针对纳米硒（SeNPs）细胞毒性评价实验成本高、可用检测数据有限的问题，本研究基于已有文献数据和结合其材料理化特性构建了关于SeNPs的细胞存活率预测模型。通过比较多种机器学习方法，我们考察了材料化学特征对细胞毒性预测性能的影响，并采用按文献来源分组的交叉验证验证模型的泛化能力。在收集到的纳米硒数据上进行实验，随机森林基础模型表现最佳，平均绝对误差（MAE）为 19.11。在此基础上，我们进一步对 9 种纳米硒材料开展人脐静脉内皮细胞（HUVEC）实验，在不同浓度和暴露时间条件下获得 351 组预测与实测配对数据，预测 MAE 为 14.51，决定系数（$R^2$）为 0.446。结果表明，该模型能够在一定程度上预测新材料的细胞存活率，在较长暴露时间下的响应变化及在不同材料和作用时间下材料间毒性差异，是进一步提高预测精度的重点。本研究将文献数据建模与新材料实验检验相结合，可为纳米硒细胞毒性预测及材料初步筛选提供参考。